\section{Individual Barnes functions and all multiple Gamma ranks}
\label{bhr:sec:barnes}

The balanced functions in the preceding reports are particular
linear combinations of Hurwitz-zeta derivatives and Bernoulli
polynomials. An individual multiple Gamma function also carries
polynomial coefficients in $x$. Theorem~\ref{bhr:thm:weighted-master}
supplies exactly those missing weights.

\subsection{Fixing the multiple Gamma normalization}

For an integer $r\ge1$, the unit-period Barnes zeta is initially
\begin{equation}\label{bhr:eq:barnes-zeta}
 \mathcal Z_r(u,x)=\sum_{n=0}^{\infty}
        \binom{n+r-1}{r-1}(n+x)^{-u},\qquad \Re u>r,\quad x>0.
\end{equation}
Define the explicit coefficient polynomials
\begin{equation}\label{bhr:eq:C-rk}
 C_{r,k}(x)=\frac1{(r-1)!}[X^k]
                      \prod_{h=1}^{r-1}(X+h-x),\qquad 0\le k<r.
\end{equation}
Expanding the binomial coefficient in powers of $n+x$ gives
\begin{equation}\label{bhr:eq:Barnes-Hurwitz}
 \mathcal Z_r(u,x)=\sum_{k=0}^{r-1}C_{r,k}(x)\zeta(u-k,x),\qquad
 L_r(x):=\partial_u\mathcal Z_r(0,x)
       =\sum_{k=0}^{r-1}C_{r,k}(x)\zeta'(-k,x).
\end{equation}
The first equality holds in the initial convergence domain and
then meromorphically. The value at zero is regular. This is the
classical Barnes/Hurwitz conversion, here written in a directly
computable polynomial form \cite{bhr:Adamchik,bhr:Vardi}.

To specify a normalization completely, put
\begin{equation}\label{bhr:eq:Q-r}
 Q_r(x)=\sum_{j=0}^{r-1}(-1)^{j+1}L_{r-j}(1)
                          \binom{x-1}{j},\qquad
 \log\Gamma_r(x)=L_r(x)+Q_r(x).
\end{equation}
We use the division convention for the multiple Gamma hierarchy:
$\Gamma_0(x)=x^{-1}$, $\Gamma_1=\Gamma$, and
$\Gamma_2=G^{-1}$. The zeta prescription
\eqref{bhr:eq:Barnes-Hurwitz}--\eqref{bhr:eq:Q-r}, not merely a functional
equation, defines the normalized functions in every higher rank.
Other standard sign conventions can be converted by inversion.

\begin{proposition}[Normalization and recurrence]\label{bhr:prop:normalization}
The functions in \eqref{bhr:eq:Q-r} satisfy
\begin{equation}\label{bhr:eq:Gamma-r-recurrence}
 \Gamma_r(1)=1,\qquad
 \Gamma_r(x+1)=\frac{\Gamma_r(x)}{\Gamma_{r-1}(x)}.
\end{equation}
Their normalization polynomials contain only rational numbers and
ordinary zeta derivatives at nonpositive integers.
\end{proposition}
\begin{proof}
Counting the multiplicities in \eqref{bhr:eq:barnes-zeta} gives
$\mathcal Z_r(u,x+1)-\mathcal Z_r(u,x)=-\mathcal Z_{r-1}(u,x)$, with
$\mathcal Z_0(u,x)=x^{-u}$. Hence
$L_r(x+1)-L_r(x)=-L_{r-1}(x)$, where $L_0=-\log x$.
The binomial difference formula gives
$Q_r(x+1)-Q_r(x)=-Q_{r-1}(x)$ for $r\ge2$.
For $r=1$, $Q_1=-\zeta'(0)=\tfrac12\log(2\pi)$ is
constant. Also $Q_r(1)=-L_r(1)$. These statements prove
\eqref{bhr:eq:Gamma-r-recurrence}, including the usual Gamma
recurrence at $r=1$. The last claim follows from
\eqref{bhr:eq:Barnes-Hurwitz} at $x=1$.
\end{proof}

The first three rows make every coefficient explicit. Here
$L=\log(2\pi)$:
\begin{align}
 L_1(x)&=\zeta'(0,x),& Q_1(x)&=L/2,\label{bhr:eq:rank1}\\
 L_2(x)&=\zeta'(-1,x)+(1-x)\zeta'(0,x),&
 Q_2(x)&=-\zeta'(-1)-\tfrac L2(x-1),\label{bhr:eq:rank2}\\
 L_3(x)&=\tfrac12\zeta'(-2,x)
       +\tfrac{3-2x}{2}\zeta'(-1,x)
       +\tfrac{(x-1)(x-2)}2\zeta'(0,x),\label{bhr:eq:rank3a}\\
 Q_3(x)&=-\tfrac12\bigl(\zeta'(-2)+\zeta'(-1)\bigr)
       +(x-1)\zeta'(-1)+\tfrac L2\binom{x-1}{2}.
       \label{bhr:eq:rank3b}
\end{align}
At rank two these equations are equivalent to the usual Vardi
identity
\begin{equation}\label{bhr:eq:Vardi}
 \log G(x)=\zeta'(-1)+(x-1)\log\Gamma(x)-\zeta'(-1,x).
\end{equation}
The familiar Barnes function has its standard asymptotic
normalization; $G(x+1)=\Gamma(x)G(x)$ and $G(1)=1$ alone
would allow multiplication by a nontrivial periodic factor.

\subsection{The all-rank resolvent theorem}

\begin{theorem}[Individual multiple Gamma transforms]\label{bhr:thm:all-Barnes}
For every integer $r\ge1$, polynomial $W$, and $z\notin\R$,
\begin{equation}\label{bhr:eq:all-Barnes}
 \boxed{\displaystyle
 \int_0^1 W(x)\log\Gamma_r(x)R_z(x)\,dx
  =-\sum_{k=0}^{r-1}
      \left.\partial_s I_{W C_{r,k}}(s,z)\right|_{s=k+1}
          +T_{WQ_r}(z).}
\end{equation}
Every term on the right is explicitly given by
\eqref{bhr:eq:weighted-master} and \eqref{bhr:eq:Tpoly}. It is a finite
combination of ordinary and depth-two polylogarithms, their first
order derivatives, Gamma factors, logarithms, and ordinary zeta
values and derivatives. All depth-two series needed at the stated
evaluation points are normally convergent.
\end{theorem}

\begin{proof}
Differentiate the ordinary integral in
Theorem~\ref{bhr:thm:weighted-master} at $s=k+1$:
\[
 \partial_s I_P(k+1,z)
       =-\int_0^1P(x)\zeta'(-k,x)R_z(x)\,dx.
\]
The shift formula at zero bounds the differentiated integrand by
a fixed polynomial in $|\log x|$ times $x$ when $k=0$, and
by an even more integrable power when $k>0$; its analytic tail
is bounded. Thus differentiation is justified locally uniformly.
Insert \eqref{bhr:eq:Barnes-Hurwitz} and \eqref{bhr:eq:Q-r}.
At every evaluation point $s=k+1\ge1$ and every polynomial
index $j\ge1$, the depth-two total order is $s+j\ge2$.
Lemma~\ref{bhr:lem:omega-normal} therefore applies with a margin,
including all the stated spectral derivatives. The sums over
$r,k,j$ are finite.
\end{proof}

This theorem answers the individual, unit-period multiple Gamma
part of source question R6. It does not assert a reduction for
arbitrary products of two multiple Gamma logarithms, nor for
arbitrary incommensurable Barnes periods.

Repeated poles are equally explicit:
\begin{equation}\label{bhr:eq:Barnes-repeated}
 \int_0^1\frac{W(x)\log\Gamma_r(x)}{(K(x)-z)^h}\,dx
   =\frac1{(h-1)!}\partial_z^{h-1}
      \int_0^1W(x)\log\Gamma_r(x)R_z(x)\,dx,
 \qquad h\ge1.
\end{equation}
All these integrals converge. Argument derivatives of
$\mathcal O_{\alpha,\beta}(q)$ can be taken directly in its
series, multiplying by powers of $n$. The formal identity
\[
 q\partial_q\mathcal O_{\alpha,\beta}
  =\mathcal O_{\alpha-1,\beta}-\mathcal O_{\alpha,\beta-1}
\]
must be kept as a paired difference when either term separately
leaves its convergence domain. The original differentiated series
still converges normally. This qualification prevents an artificial
divergence in implementations at total order two.

\subsection{A compact formula for the ordinary Barnes function}

Only one depth-two derivative is needed at rank two. Define
\begin{equation}\label{bhr:eq:Dq}
 \mathcal D(q)=
    \left.(\partial_\alpha+\partial_\beta)
                      \mathcal O_{\alpha,\beta}(q)\right|_{\alpha=\beta=1}
  =-\sum_{n,m\ge1}\frac{q^n\log(m(m+n))}{m(m+n)}.
\end{equation}
It is holomorphic for $|q|<1$ and vanishes at zero. In the next
formulas all polylogarithms have argument $q$. Set
\begin{align}
 \mathcal A(q)
  &=\frac{\mathcal D(q)-\Li_2^{[1]}(q)
       -(\gamma+L+1)\Li_2(q)-\Li_1(q)\Li_1^{[1]}(q)}{2\pi^2}
           -\frac14\Li_1(q),\label{bhr:eq:Aq}\\
 \mathcal S(q)
  &=\frac{\Li_1^{[1]}(q)-(\gamma+2L)\Li_1(q)
            -\tfrac12\bigl(\Li_2(q)+\Li_1(q)^2\bigr)}{2\pi},
        \label{bhr:eq:Sq}\\
 \mu_G&=2\zeta'(-1)-\frac1{12}-\frac L4.
       \label{bhr:eq:meanG}
\end{align}

\begin{theorem}[The individual $G$ resolvent]\label{bhr:thm:G-resolvent}
With the coordinates \eqref{bhr:eq:coordinates},
\begin{equation}\label{bhr:eq:G-resolvent}
 \boxed{\displaystyle
 \int_0^1\frac{\log G(x)}{\pi\cot(\pi x)-z}\,dx
    =-\frac{\mu_G}{a}
       +\frac c{2q}\bigl(\mathcal A(q)+\epsilon\ii\mathcal S(q)\bigr).}
\end{equation}
The quotient at $q=0$ is removable. The formula is holomorphic
on each open half-plane and uses the real logarithm of $G(x)$
on the integration interval.
\end{theorem}

\begin{proof}
We give a summation proof that also connects the formula to the
known Barnes Fourier expansion \cite[Theorems 2--3]{bhr:Mezo}.
Write
\[
 \log G(x)=\mu_G+\sum_{n\ge1}
                  \bigl(a_n\cos(2\pi nx)+b_n\sin(2\pi nx)\bigr).
\]
Its coefficients are
\begin{align}
 a_n&=\frac{\frac12\log n-(\gamma+L)-1}{2\pi^2n^2}
                -\frac1{4n}-\frac{T_n}{\pi^2},\label{bhr:eq:anG}\\
 b_n&=\frac{\frac1{2n}-\gamma-2L-\log n-H_n}{2\pi n},
          \label{bhr:eq:bnG}\\
 T_n&=\sum_{\substack{m\ge1\\m\ne n}}
                          \frac{\log m}{m^2-n^2}.
          \label{bhr:eq:Tn}
\end{align}
The $m$-sum in $T_n$ converges absolutely. These coefficients
can also be recovered from Theorem~\ref{bhr:thm:weighted-master}
by differentiating the $P=x-1$ formula at $s=1$, the $P=1$
formula at $s=2$, and using \eqref{bhr:eq:Vardi}. Thus their use
does not assume a new Barnes Fourier theorem.

To evaluate their generating functions, differentiate the
$j=1$ case of \eqref{bhr:eq:lattice-splitting} and the definition
of $\mathcal O_{1,s}$. Partial fractions give the exact identity
\begin{equation}\label{bhr:eq:T-generator}
 \sum_{n\ge1}q^nT_n
   =\frac12\Li_1(q)\Li_1^{[1]}(q)
       -\frac12\mathcal D(q)+\frac14\Li_2^{[1]}(q).
\end{equation}
For clarity, the scalar calculation behind it is
\[
 \frac1{m^2-n^2}=
     \frac1{2m}\left(\frac1{m-n}+\frac1{m+n}\right).
\]
The excluded term $m=n$ in the second sum is
$\log n/(2n^2)$ before multiplication by $1/2$; it supplies
the last term in \eqref{bhr:eq:T-generator}. This diagonal term
is necessary for the coefficient of $\Li_2^{[1]}$.

Next,
\begin{equation}\label{bhr:eq:harmonic-generator}
 \sum_{n\ge1}\frac{H_nq^n}{n}
      =\Li_2(q)+\frac12\Li_1(q)^2.
\end{equation}
It follows either by integrating
$\sum_{n\ge1}H_nq^n=\Li_1(q)/(1-q)$ or by the depth-two
stuffle identity. Equations \eqref{bhr:eq:anG}--\eqref{bhr:eq:harmonic-generator}
give $\sum a_nq^n=\mathcal A(q)$ and
$\sum b_nq^n=\mathcal S(q)$.

The function $\log G$ is integrable and square integrable:
near zero it equals $\log x+O(x)$, while it is analytic at one.
Pair its Fourier coefficients with the uniformly convergent
geometric series \eqref{bhr:eq:resolvent-fourier}. This yields
\eqref{bhr:eq:G-resolvent}. The constant term is
\eqref{bhr:eq:meanG}; for example integrate \eqref{bhr:eq:Vardi}, use
$\int_0^1\zeta'(-1,x)\,dx=0$ and
\[
 \int_0^1x\log\Gamma(x)\,dx
       =\frac L4+\zeta'(-1)-\frac1{12}.
\]
The latter follows by integrating the Hurwitz argument ladder.
All generating functions used above are $O(q)$, so the quotient
at zero is removable. Normal convergence proves holomorphy.
\end{proof}

\subsection{An ordinary-polylogarithm real component}

\begin{corollary}[A one-parameter exact real integral]\label{bhr:cor:G-real}
For $b>0$, let $q=(b-\pi)/(b+\pi)\in(-1,1)$. Then
\begin{equation}\label{bhr:eq:G-real}
 \boxed{\displaystyle
 \int_0^1\frac{K(x)\log G(x)}{K(x)^2+b^2}\,dx
 =\frac{\Li_1^{[1]}(q)-(\gamma+2L)\Li_1(q)
       -\frac12\bigl(\Li_2(q)+\Li_1(q)^2\bigr)}
               {2q(b+\pi)^2}.}
\end{equation}
At $b=\pi$ the right side means
$-(\gamma+2L+1/2)/(8\pi^2)$. In particular,
\begin{equation}\label{bhr:eq:G-threepi}
 \boxed{\displaystyle
 \int_0^1\frac{\pi\cot(\pi x)\log G(x)}
                      {\pi^2\cot^2(\pi x)+9\pi^2}\,dx
 =\frac{\Li_1^{[1]}(1/2)-(\gamma+2\log(2\pi))\log2
                    -\pi^2/24-(\log2)^2/4}{16\pi^2}.}
\end{equation}
\end{corollary}
\begin{proof}
At $z=\ii b$, $a=\ii(b+\pi)$ and
$c=-2\pi\ii/(b+\pi)^2$. The functions $\mathcal A,
\mathcal S$ are real at real $q$. Taking the real part of
\eqref{bhr:eq:G-resolvent} therefore keeps only
$\pi\mathcal S(q)/(q(b+\pi)^2)$.
Use $\Li_1(q)/q\to1$, $\Li_2(q)/q\to1$ and
$\Li_1^{[1]}(q)/q\to0$ at zero. For $b=3\pi$, use
$q=1/2$, $\Li_1(1/2)=\log2$ and
$\Li_2(1/2)=\pi^2/12-(\log2)^2/2$.
\end{proof}

The integral \eqref{bhr:eq:G-threepi} is approximately
$-0.02312768479852353748486907193849$.
The complete complex transform at $z=3\pi\ii$ is approximately
\[
 -0.02312768479852353748486907193849
   -0.05510836371472001971368271205301\,\ii.
\]
These decimals are diagnostics, not definitions or exact-value
certificates. The symmetry of $K$ under $x\mapsto1-x$ explains
why its real component isolates the simpler sine coefficients.

\subsection{Normalized antiderivatives remain explicit}

The standard balanced primitives are
\begin{equation}\label{bhr:eq:balanced-B}
 \mathscr B_k(x)=\frac{\zeta'(-k,x)+H_k\zeta(-k,x)}{k!},
 \qquad H_0=0.
\end{equation}
They satisfy $\mathscr B_0=\log\Gamma-L/2$ and
$\mathscr B_k'=\mathscr B_{k-1}$ for $k\ge1$
\cite{bhr:EMIntegrals,bhr:EMGeneralized}. For a polynomial $P$, set
\begin{equation}\label{bhr:eq:weighted-antiderivative}
 \mathscr A_{P,k}(x)=
    \sum_{j=0}^{\deg P}(-1)^jP^{(j)}(x)\mathscr B_{k+j+1}(x).
\end{equation}
Termwise differentiation telescopes, so
$\mathscr A_{P,k}'=P\mathscr B_k$.
Consequently a primitive of $P(x)\zeta'(-k,x)$, normalized
to vanish at $x=1$, is
\begin{align}\label{bhr:eq:normalized-zeta-primitive}
 k!\bigl(\mathscr A_{P,k}(x)-\mathscr A_{P,k}(1)\bigr)
    -H_k\int_1^xP(t)\zeta(-k,t)\,dt.
\end{align}
The remaining integral is polynomial because
$\zeta(-k,t)=-B_{k+1}(t)/(k+1)$.
Apply \eqref{bhr:eq:normalized-zeta-primitive} to the finitely many
terms in \eqref{bhr:eq:Barnes-Hurwitz}, and integrate $Q_r$ as a
polynomial. This gives a fully normalized primitive of every
$\log\Gamma_r$. Repetition gives every iterated primitive,
without unspecified additive constants. Their resolvent
transforms again follow from Theorem~\ref{bhr:thm:all-Barnes} and
Theorem~\ref{bhr:thm:weighted-master}. The primitive operation is
classical integration by parts; the point here is its compatibility
with the explicit individual-function normalization.
